\begin{exercise}{105}
Show that for any two simplicial sets, $X\star Y$ is a retract of
$X\diamond Y$.\footnote{Replace the original ``Show that'' by ``Disprove that'':
for $X=\Lambda^2_1$ and $Y=\Delta^0$, no map $X\star Y\to X\diamond Y$
is injective on vertices, as proved below, so no retraction exists.}

\begin{proof}
Take $X=\Lambda^2_1$ and $Y=\Delta^0$, and label the cone vertex by $3$.
Then
\[
 A=X\star\Delta^0
 =\Delta^{\{0,1,3\}}\cup_{\Delta^{\{1,3\}}}\Delta^{\{1,2,3\}}
 \subseteq\Delta^3.
\]
By Definition~2.5.1,
\[
 B=X\diamond\Delta^0
 =(X\times\Delta^1)\amalg_{X\times\{1\}}\Delta^0.
\]
Thus an $n$-simplex of $B$ is represented by $[u,t]$, where
$u\in X_n$ and $t\in(\Delta^1)_n$.
All pairs with $t=(1,\ldots,1)$ are identified, and there are no other
identifications. The vertices are $0,1,2,3$, with $3$ the collapsed end.
Recall that simplices of $X$ are nondecreasing tuples in $\{0,1,2\}$
which do not contain both $0$ and $2$.

Every edge of $B$ with distinct endpoints points forward in the order
$0<1<2<3$.
If a map $i\colon A\to B$ is injective on vertices, the three edges
$01,12,23$ of $A$ therefore force
$i(0)<i(1)<i(2)<i(3)$, so $i$ fixes every vertex.

The only $2$-simplices of $B$ with respective vertex sequences
$(0,1,3)$ and $(1,2,3)$ are
\[
 S=[(0,1,1),(0,0,1)],\qquad
 T=[(1,2,2),(0,0,1)].
\]
Indeed, the second tuple must be $(0,0,1)$, and the restrictions on
tuples in $X$ force the displayed first tuples.
Hence $i$ must send the two triangles of $A$ to $S$ and $T$.
Their common edge $13$ would then have both images
\[
 d_0S=[(1,1),(0,1)]\ne[(1,2),(0,1)]=d_1T,
\]
which are distinct by the quotient description of $B$.
This is impossible. Since the section of any retraction is injective
on vertices, $A$ cannot be a retract of $B$.
\end{proof}
\end{exercise}
